\chapter{Редуцированные состояния, энтропия запутанности и унитарная эволюция}

Матрица плотности впервые появляется как естественное описание неполной подготовки: квантовая система подготовлена, но мы не знаем в точности, как именно. Запутанность придаёт тому же математическому объекту более глубокую роль. Даже когда состояние составной системы известно полностью, подсистема, вообще говоря, не имеет собственного вектора состояния. Её редуцированная матрица плотности содержит всё, что можно предсказать локально. Мы выведем это утверждение покомпонентно, проверим его на двух спиновых состояниях и используем энтропию, чтобы отличать состояния-произведения от запутанных. Затем мы перейдём ко второму вопросу: как квантовое состояние изменяется со временем. Сохранение скалярных произведений приведёт к унитарности, эрмитову генератору и уравнению Schr\"odinger.

\section{Матрицы плотности и неполная подготовка}

Предположим, что электрон помещён в сильное магнитное поле и ему дали выровняться вдоль него. Направление поля выбрал и электрон подготовил кто-то другой, но нам это направление не сообщил. Поэтому спин может находиться в совершенно определённом чистом состоянии, хотя мы не знаем, какое именно это состояние. Мы можем знать распределение вероятностей по возможным способам подготовки, а можем не знать вообще ничего. Матрица плотности \(\rho\) — это квантовое описание именно такого рода неполной информации.

Классический аналог — это плотность вероятности на фазовом пространстве или распределение вероятностей по дискретным конфигурациям. Квантовый аналог — оператор с тремя существенными свойствами:
\begin{equation}
\rho^\dagger=\rho,
\qquad
\operatorname{Tr}\rho=1,
\qquad
\lambda_i\geq 0,
\end{equation}
где \(\lambda_i\) — собственные значения \(\rho\). Эрмитовость делает собственные значения вещественными и задаёт ортонормированный собственный базис. Положительность делает их допустимыми вероятностями, а нормировка даёт
\begin{equation}
\sum_{i=1}^{n}\lambda_i=1.
\end{equation}
В собственном базисе \(\rho\) величина \(\lambda_i\) есть вероятность, приписанная \(i\)-му собственному состоянию.

Полное незнание — это случай, когда ни одному направлению в гильбертовом пространстве не отдаётся предпочтение:
\begin{equation}
\rho_{\mathrm{max}}=\frac{1}{n}I.
\end{equation}
Каждое собственное значение равно \(1/n\). На противоположном пределе полное знание представляется нормированным состоянием \(\ket{\psi}\), причём
\begin{equation}
\rho_\psi=\ket{\psi}\bra{\psi}.
\end{equation}
Это проектор на одномерное подпространство, порождённое \(\ket{\psi}\).

\begin{classroomqa}
\lecturetimestamp{00:06:47}
\audiencequestion{Можно ли описать подготовку как выбор пятьдесят на пятьдесят между двумя состояниями, которые не ортогональны, например между двумя направлениями спина, разделёнными углом \(45^\circ\)?}
\lecturerresponse{Такая процедура подготовки возможна, но эти два веса, вообще говоря, не являются собственными значениями результирующей матрицы плотности. Сначала строят взвешенную сумму проекторов, а затем диагонализуют её. Прямая вероятностная интерпретация собственных значений относится к ортонормированному собственному базису \(\rho\), а не обязательно к неортогональным состояниям, указанным тем, кто выполнял подготовку.}
\end{classroomqa}

\begin{classroomqa}
\lecturetimestamp{00:07:47}
\audiencequestion{Чем чистое состояние отличается от смешанного?}
\lecturerresponse{У матрицы плотности чистого состояния одно ненулевое собственное значение, равное \(1\). У смешанного состояния ненулевых собственных значений больше одного. В старой литературе использовалось немецкое слово \emph{Gemisch}; современный термин — просто \emph{смешанное состояние}.}
\end{classroomqa}

Проверим спектр проектора на чистое состояние. Нормировка даёт
\begin{equation}
\rho_\psi\ket{\psi}
=\ket{\psi}\braket{\psi}{\psi}
=\ket{\psi},
\end{equation}
поэтому \(\ket{\psi}\) — собственный вектор с собственным значением \(1\). Если \(\ket{\phi}\) ортогонален \(\ket{\psi}\), то
\begin{equation}
\rho_\psi\ket{\phi}
=\ket{\psi}\braket{\psi}{\phi}
=0.
\end{equation}
Следовательно, любое направление, ортогональное \(\ket{\psi}\), имеет нулевое собственное значение.

\begin{figure}[htbp]
\centering
\includegraphics[width=0.72\textwidth]{lecture_08_figure_06.png}
\caption{Матрица плотности чистого состояния действует как проектор: \(\rho\ket{\psi}=\ket{\psi}\).}
\lectureframe{00:09:32}
\end{figure}

\section{Средние значения через след}

Рабочее правило для матрицы плотности имеет вид
\begin{equation}
\overline{M}=\operatorname{Tr}(\rho M),
\end{equation}
где \(M\) — наблюдаемая. Для двух множителей циклическое свойство следа даёт
\begin{equation}
\operatorname{Tr}(\rho M)=\operatorname{Tr}(M\rho),
\end{equation}
независимо от того, коммутируют \(\rho\) и \(M\) или нет.

В любом ортонормированном базисе \(\{\ket{i}\}\),
\begin{equation}
\operatorname{Tr}X=\sum_i\bra{i}X\ket{i}.
\end{equation}
Численное значение не зависит от выбора базиса. Для чистого состояния
\begin{align}
\overline{M}
&=\operatorname{Tr}\!\left(\ket{\psi}\bra{\psi}M\right) \\
&=\sum_i\braket{i}{\psi}\bra{\psi}M\ket{i} \\
&=\bra{\psi}M\left(\sum_i\ket{i}\bra{i}\right)\ket{\psi} \\
&=\bra{\psi}M\ket{\psi}.
\end{align}
Разложение единицы
\begin{equation}
\sum_i\ket{i}\bra{i}=I,
\end{equation}
сводит правило с матрицей плотности обратно к привычному среднему значению для чистого состояния.

\begin{figure}[htbp]
\centering
\includegraphics[width=0.74\textwidth]{lecture_08_figure_07.png}
\caption{Вычисление следа для проектора на чистое состояние и разложение единицы.}
\lectureframe{00:14:25}
\end{figure}

\begin{classroomqa}
\lecturetimestamp{00:14:38}
\audiencequestion{Должна ли величина слева быть \(\overline{M}\), то есть средним значением \(M\)?}
\lecturerresponse{Да. След — это не сам оператор \(M\); это скалярное среднее значение \(\overline{M}\).}
\end{classroomqa}

Для произвольной матрицы плотности выберем её собственный базис:
\begin{equation}
\rho\ket{j}=\lambda_j\ket{j}.
\end{equation}
Вставим полную систему состояний между \(\rho\) и \(M\):
\begin{align}
\overline{M}
&=\sum_i\bra{i}\rho M\ket{i} \\
&=\sum_{i,j}\bra{i}\rho\ket{j}\bra{j}M\ket{i} \\
&=\sum_{i,j}\lambda_j\delta_{ij}\bra{j}M\ket{i} \\
&=\sum_j\lambda_j\bra{j}M\ket{j}.
\end{align}
Это обычное взвешенное среднее: \(\lambda_j\) — вероятность \(j\)-го собственного состояния, а \(\bra{j}M\ket{j}\) — среднее значение \(M\) в этом состоянии.

\begin{classroomqa}
\lecturetimestamp{00:18:12}
\audiencequestion{Почему обе метки в последнем матричном элементе становятся \(j\), и по какому индексу всё ещё ведётся суммирование?}
\lecturerresponse{Поскольку \(\rho\) диагональна в этом базисе, \(\bra{i}\rho\ket{j}=\lambda_j\delta_{ij}\), выживают только члены с \(i=j\). Исходная двойная сумма сводится к одной сумме, которую можно называть суммой по \(j\). Название немого индекса суммирования не имеет физического значения.}
\end{classroomqa}

\section{Энтропия как недостающая информация}

Для классического распределения вероятностей \(\{p_i\}\) определим
\begin{equation}
S=-\sum_i p_i\log p_i,
\end{equation}
где \(0\log 0\) понимается как его предельное значение \(0\). Если одна вероятность равна \(1\), а все остальные обращаются в нуль, то \(S=0\): состояние известно. Если все \(n\) возможностей равновероятны,
\begin{equation}
p_i=\frac{1}{n},
\qquad
S=-n\left(\frac{1}{n}\log\frac{1}{n}\right)=\log n.
\end{equation}
Таким образом, энтропия растёт по мере того, как распределение вероятностей размазывается по большему числу значимых возможностей. Грубо говоря, \(e^S\) — это эффективное число состояний с сопоставимыми весами, когда используются натуральные логарифмы.

\begin{figure}[htbp]
\centering
\includegraphics[width=0.68\textwidth]{lecture_08_figure_08.png}
\caption{Энтропия для распределения вероятностей и значение \(S=\log n\) при равных вероятностях.}
\lectureframe{00:21:08}
\end{figure}

Для матрицы плотности то же определение применяется к её собственным значениям:
\begin{equation}
S(\rho)=-\sum_i\lambda_i\log\lambda_i
=-\operatorname{Tr}(\rho\log\rho).
\end{equation}
Первое выражение нам непосредственно понадобится; форма через след подчёркивает, что ответ не зависит от выбора базиса.

\section{Составная система и локальный эксперимент}

Рассмотрим две подсистемы, \(A\) и \(B\). Базисный вектор записывается как \(\ket{a,b}\). Если у \(A\) есть \(N\) базисных состояний, а у \(B\) есть \(n\), то у составной системы есть \(Nn\) базисных состояний. Эти две системы не обязаны быть одинаковыми: одна может быть спином с двумя состояниями, а другая — атомом с тремя состояниями.

Наиболее общее чистое состояние пары имеет вид
\begin{equation}
\ket{\Psi}
=\sum_{a,b}\psi(a,b)\ket{a,b},
\qquad
\sum_{a,b}|\psi(a,b)|^2=1.
\end{equation}
Пусть теперь Алиса измеряет наблюдаемую \(M_A\), действующую только на подсистему \(A\). В составном гильбертовом пространстве этот оператор имеет вид \(M_A\otimes I_B\). Его матричные элементы равны
\begin{equation}
\bra{a',b'}M_A\otimes I_B\ket{a,b}
=\bra{a'}M_A\ket{a}\,\delta_{b'b}.
\end{equation}
Индекс \(B\) пассивен.

Выпишем среднее значение полностью:
\begin{align}
\expect{M_A}
&=\bra{\Psi}M_A\otimes I_B\ket{\Psi} \\
&=\sum_{a',b'}\sum_{a,b}
\psi^*(a',b')\psi(a,b)
\bra{a'}M_A\ket{a}\delta_{b'b} \\
&=\sum_{a',a,b}
\psi^*(a',b)\psi(a,b)
\bra{a'}M_A\ket{a}.
\end{align}

\begin{classroomqa}
\lecturetimestamp{00:29:12}
\audiencequestion{Почему метки у бра стоят со штрихами, а у кета — нет?}
\lecturerresponse{Бра и кет до вычисления матричного элемента входят в независимые суммы, поэтому для них нужны отдельные немые метки. Штрихи отличают индексы бра от индексов кета. Затем оператор и скалярные произведения определяют, какие из этих меток должны совпасть.}
\end{classroomqa}

Зафиксируем \(a\) и \(a'\) и выполним суммирование по \(b\). Результат — матрица на подсистеме \(A\):
\begin{equation}
\boxed{
\rho_A(a,a')
=\sum_b\psi(a,b)\psi^*(a',b)
}.
\end{equation}
Это редуцированная матрица плотности \(A\) или, что то же самое, частичный след проектора на чистое состояние по \(B\):
\begin{equation}
\rho_A=\operatorname{Tr}_B\ket{\Psi}\bra{\Psi}.
\end{equation}
Локальное среднее значение принимает вид
\begin{equation}
\expect{M_A}
=\sum_{a,a'}\rho_A(a,a')\bra{a'}M_A\ket{a}
=\operatorname{Tr}_A(\rho_A M_A).
\end{equation}

\begin{figure}[htbp]
\centering
\includegraphics[width=0.76\textwidth]{lecture_08_figure_09.png}
\caption{Индекс \(B\) суммируется, оставляя матрицу по двум индексам \(A\).}
\lectureframe{00:33:50}
\end{figure}

\begin{classroomqa}
\lecturetimestamp{00:35:39}
\audiencequestion{Является ли локальная наблюдаемая просто единичным оператором в подсистеме \(B\)?}
\lecturerresponse{Именно. \(M_A\otimes I_B\) может изменять координату \(A\), но он не изменяет \(b\) и вообще от неё не зависит. Поэтому его матричный элемент в пространстве \(B\) равен \(\delta_{b'b}\).}
\end{classroomqa}

Все предсказания для экспериментов, ограниченных \(A\), поэтому содержатся в \(\rho_A\). Бобу не нужно измерять свою подсистему. Уже одно неиспользование степеней свободы \(B\) вынуждает описание Алисы принять форму матрицы плотности.

\section{Чистое целое может иметь смешанные части}

Составная система находится в чистом состоянии \(\ket{\Psi}\), однако у \(\rho_A\), вообще говоря, будет несколько ненулевых собственных значений. Полное знание целого не влечёт существования вектора состояния для каждой части. У кошки и ружья может быть один общий запутанный вектор состояния; тогда одна только кошка описывается редуцированной матрицей плотности.

\begin{classroomqa}
\lecturetimestamp{00:38:25}
\audiencequestion{Можно ли заменить \(\operatorname{Tr}(M\rho)\) на \((\operatorname{Tr}M)(\operatorname{Tr}\rho)\)?}
\lecturerresponse{Нет. След произведения, вообще говоря, не равен произведению следов. Цикличность допускает \(\operatorname{Tr}(M\rho)=\operatorname{Tr}(\rho M)\), но не позволяет разложить след на множители.}
\end{classroomqa}

\begin{classroomqa}
\lecturetimestamp{00:38:49}
\audiencequestion{Начинали ли мы только с базисного состояния \(\ket{a,b}\) или с наиболее общего чистого состояния?}
\lecturerresponse{С наиболее общего чистого состояния составной системы: \(\ket{\Psi}=\sum_{a,b}\psi(a,b)\ket{a,b}\). Одно \(\ket{a,b}\) — это лишь один базисный вектор. Неожиданно то, что даже этот полностью заданный составной вектор может давать смешанное редуцированное состояние.}
\end{classroomqa}

Контраст с классической механикой резок. Полная классическая информация о составной системе фиксирует каждую её координату, а значит, и каждую подсистему. От классического распределения вероятностей, конечно, можно перейти к маргинальному, но если исходное распределение сосредоточено в одной точке фазового пространства, то и любое маргинальное распределение сосредоточено при одном определённом значении. Квантовые чистые состояния не обязаны обладать этим свойством.

\begin{classroomqa}
\lecturetimestamp{00:40:46}
\audiencequestion{Чем суммирование по \(B\) отличается от интегрирования классического распределения по координате--скорости по скорости?}
\lecturerresponse{Операции выглядят похожими, но в пределе чистого состояния они различаются. Полностью известная классическая точка фазового пространства фиксирует и координату, и скорость, так что любая маргиналь остаётся определённой. Полностью известное запутанное квантовое состояние может давать нетривиальное смешанное состояние для любой подсистемы. Частичный след тем самым показывает, что в квантовой теории не действует классическое следствие: из чистоты целого не вытекает чистота частей.}
\end{classroomqa}

\section{Состояния-произведения: исключительный случай}

Когда \(\rho_A\) остаётся чистым? Ответ таков: волновая функция должна факторизоваться:
\begin{equation}
\psi(a,b)=\phi(a)\chi(b).
\end{equation}
Предположим, что \(\phi\) и \(\chi\) по отдельности нормированы. Подстановка в редуцированную матрицу плотности даёт
\begin{align}
\rho_A(a,a')
&=\sum_b\phi(a)\chi(b)\phi^*(a')\chi^*(b) \\
&=\phi(a)\phi^*(a')\sum_b|\chi(b)|^2 \\
&=\phi(a)\phi^*(a').
\end{align}
Следовательно,
\begin{equation}
\rho_A=\ket{\phi}\bra{\phi}.
\end{equation}
У неё одно собственное значение \(1\), все остальные собственные значения равны \(0\), и энтропия равна нулю.

\begin{figure}[htbp]
\centering
\includegraphics[width=0.74\textwidth]{lecture_08_figure_10.png}
\caption{Для факторизованной волновой функции сумма по \(B\) сводится к нормировке \(\chi\).}
\lectureframe{00:46:35}
\end{figure}

\begin{classroomqa}
\lecturetimestamp{00:44:17}
\audiencequestion{Является ли \(\phi(a)\) само по себе состоянием системы \(A\)?}
\lecturerresponse{Это компонент волновой функции этого состояния. Более явно, \(\ket{\phi}=\sum_a\phi(a)\ket{a}\), и аналогично \(\ket{\chi}=\sum_b\chi(b)\ket{b}\). Факторизация означает \(\ket{\Psi}=\ket{\phi}\otimes\ket{\chi}\), то есть состояние-произведение без запутанности между \(A\) и \(B\).}
\end{classroomqa}

Для любой локальной наблюдаемой величины,
\begin{equation}
\expect{M_A}
=\operatorname{Tr}(\rho_A M_A)
=\bra{\phi}M_A\ket{\phi}.
\end{equation}
Подсистема \(A\) ведёт себя точно так, как если бы она находилась в чистом состоянии \(\ket{\phi}\). Для чистого двухчастичного состояния энтропия \(\rho_A\), следовательно, служит тестом на структуру произведения: нуль означает состояние-произведение, а ненулевая энтропия означает запутанность. Значения, близкие к нулю, указывают на слабую запутанность; значения, близкие к максимуму, — на состояние, далёкое от любого произведения относительно этого разбиения на две части.

\section{Синглет: чистое целое и максимальное локальное незнание}

Возьмём двухспиновый синглет
\begin{equation}
\ket{\Psi^-}
=\frac{1}{\sqrt{2}}
\left(
\ket{\uparrow\downarrow}
-\ket{\downarrow\uparrow}
\right).
\end{equation}
Его ненулевые компоненты волновой функции равны
\begin{equation}
\psi(\uparrow,\downarrow)=\frac{1}{\sqrt{2}},
\qquad
\psi(\downarrow,\uparrow)=-\frac{1}{\sqrt{2}},
\end{equation}
тогда как \(\psi(\uparrow,\uparrow)=\psi(\downarrow,\downarrow)=0\).

Выполним частичный след по спину \(2\). Диагональные элементы равны
\begin{align}
\rho_1(\uparrow,\uparrow)
&=|\psi(\uparrow,\uparrow)|^2
 +|\psi(\uparrow,\downarrow)|^2
=\frac{1}{2}, \\
\rho_1(\downarrow,\downarrow)
&=|\psi(\downarrow,\uparrow)|^2
 +|\psi(\downarrow,\downarrow)|^2
=\frac{1}{2}.
\end{align}
Внедиагональный элемент равен
\begin{align}
\rho_1(\uparrow,\downarrow)
&=\psi(\uparrow,\uparrow)\psi^*(\downarrow,\uparrow)
 +\psi(\uparrow,\downarrow)\psi^*(\downarrow,\downarrow) \\
&=0,
\end{align}
и его комплексно-сопряжённый тоже обращается в нуль. Следовательно,
\begin{equation}
\rho_1
=
\begin{pmatrix}
\frac12&0\\[2pt]
0&\frac12
\end{pmatrix}
=\frac12 I.
\end{equation}

\begin{figure}[htbp]
\centering
\includegraphics[width=0.76\textwidth]{lecture_08_figure_11.png}
\caption{Четыре элемента редуцированной матрицы плотности одного спина в синглетном состоянии.}
\lectureframe{00:55:35}
\end{figure}

\begin{classroomqa}
\lecturetimestamp{00:55:58}
\audiencequestion{Какова энтропия редуцированного состояния одного спина?}
\lecturerresponse{Его два собственных значения оба равны \(1/2\), поэтому
\[
S_1=-2\left(\frac12\log\frac12\right)=\log 2.
\]
Это максимальная энтропия, возможная для двухсостоянийной подсистемы. Максимально смешанное состояние характеризуется равными \emph{собственными значениями}, а не равными матричными элементами в произвольном базисе.}
\end{classroomqa}

Среднее значение спина обращается в нуль вдоль любой оси:
\begin{equation}
\expect{\bm{\sigma}\cdot\hat{\bm n}}
=\operatorname{Tr}\!\left(
\rho_1\,\bm{\sigma}\cdot\hat{\bm n}
\right)
=\frac12\operatorname{Tr}
\left(\bm{\sigma}\cdot\hat{\bm n}\right)
=0.
\end{equation}

\begin{classroomqa}
\lecturetimestamp{00:58:37}
\audiencequestion{Почему след \(\bm{\sigma}\cdot\hat{\bm n}\) равен нулю для любого направления?}
\lecturerresponse{Каждая матрица Паули имеет нулевой след. Поскольку \(\bm{\sigma}\cdot\hat{\bm n}=n_1\sigma_1+n_2\sigma_2+n_3\sigma_3\), из линейности следа следует нуль для любого единичного вектора \(\hat{\bm n}\). Поэтому спин с равной вероятностью обнаруживается параллельным или антипараллельным любой выбранной оси.}
\end{classroomqa}

Вся пара имеет определённое чистое состояние и нулевую энтропию. Каждый спин сам по себе максимально смешан. Это не незнание, вызванное дефектной записью о приготовлении; это локальное состояние, обусловленное запутанностью.

\section{Состояние-произведение с равными амплитудами}

Теперь рассмотрим
\begin{equation}
\ket{\Phi}
=\frac12\left(
\ket{\uparrow\uparrow}
+\ket{\uparrow\downarrow}
+\ket{\downarrow\uparrow}
+\ket{\downarrow\downarrow}
\right).
\end{equation}
Все четыре амплитуды равны \(1/2\), поэтому нормировка следует из четырёх вкладов по \(1/4\). Частичный след по спину \(2\) даёт
\begin{equation}
\rho_1
=
\begin{pmatrix}
\frac12&\frac12\\[2pt]
\frac12&\frac12
\end{pmatrix}.
\end{equation}
Её след равен \(1\), но
\begin{equation}
\det\rho_1
=\frac14-\frac14
=0.
\end{equation}

\begin{classroomqa}
\lecturetimestamp{01:02:23}
\audiencequestion{Как здесь след и определитель определяют два собственных значения?}
\lecturerresponse{Для матрицы \(2\times2\) сумма собственных значений равна следу, а их произведение — определителю. Следовательно, \(\lambda_1+\lambda_2=1\) и \(\lambda_1\lambda_2=0\), что вынуждает пару \(\{1,0\}\). Редуцированное состояние чистое, и его энтропия равна нулю.}
\end{classroomqa}

\begin{classroomqa}
\lecturetimestamp{01:03:18}
\audiencequestion{Как на самом деле ориентированы два незапутанных спина?}
\lecturerresponse{Оба они направлены вдоль положительной оси \(x\). Это можно увидеть из факторизации,
\[
\ket{\Phi}
=\frac{\ket{\uparrow}+\ket{\downarrow}}{\sqrt2}
\otimes
\frac{\ket{\uparrow}+\ket{\downarrow}}{\sqrt2},
\]
или установить экспериментально, вычислив средние значения матриц Паули.}
\end{classroomqa}

Для \(z\)-компоненты,
\begin{equation}
\sigma_3=
\begin{pmatrix}
1&0\\[2pt]
0&-1
\end{pmatrix},
\qquad
\operatorname{Tr}(\rho_1\sigma_3)=0.
\end{equation}
Для \(x\)-компоненты,
\begin{equation}
\sigma_1=
\begin{pmatrix}
0&1\\[2pt]
1&0
\end{pmatrix},
\end{equation}
и
\begin{align}
\rho_1\sigma_1
&=
\begin{pmatrix}
\frac12&\frac12\\[2pt]
\frac12&\frac12
\end{pmatrix}
\begin{pmatrix}
0&1\\[2pt]
1&0
\end{pmatrix}
=
\begin{pmatrix}
\frac12&\frac12\\[2pt]
\frac12&\frac12
\end{pmatrix}, \\
\operatorname{Tr}(\rho_1\sigma_1)&=1.
\end{align}

\begin{figure}[htbp]
\centering
\includegraphics[width=0.70\textwidth]{lecture_08_figure_03.png}
\caption{Вычисление с матрицей \(\sigma_1\) для редуцированного состояния с равными амплитудами.}
\lectureframe{01:05:13}
\end{figure}

Таким образом, \(\expect{\sigma_1}=1\), тогда как \(\expect{\sigma_2}=\expect{\sigma_3}=0\). Отклонение амплитуд от точной факторизации, вообще говоря, меняет собственные значения с \(\{1,0\}\) на два ненулевых значения. Возникающая малая энтропия измеряет малую степень запутанности.

\section{Чья это энтропия?}

Энтропия, только что вычисленная, относится к подсистеме. В примере с синглетом,
\begin{equation}
S(\rho_{12})=0,
\qquad
S(\rho_1)=S(\rho_2)=\log 2.
\end{equation}
Целое чисто, тогда как каждая часть смешана. Поэтому энтропия запутанности не подчиняется наивному аддитивному правилу \(S_{12}=S_1+S_2\). Для чистого двухчастичного состояния две редуцированные матрицы плотности имеют одинаковые ненулевые собственные значения, поэтому их энтропии равны, хотя энтропия целого обращается в нуль.

\begin{classroomqa}
\lecturetimestamp{01:09:03}
\audiencequestion{Является ли любая матрица плотности в конечном счёте матрицей плотности запутанности?}
\lecturerresponse{С операциональной точки зрения можно описывать обычную неопределённость, не отслеживая носитель недостающей записи. Если включить аппарат приготовления, окружение или наблюдателя, эту неопределённость часто можно рассматривать как запутанность с ненаблюдаемыми степенями свободы. Математически эта идея точна: любая матрица плотности допускает пурификацию в большем гильбертовом пространстве, хотя такая пурификация не определяет однозначно физическую историю приготовления.}
\end{classroomqa}

\section{От дискретного классического движения к квантовому движению}

Теперь мы временно оставим в стороне запутанность и спросим, как состояния меняются со временем. Для классической системы с конечным числом конфигураций обратимая эволюция есть перестановка. У каждого состояния есть одно следующее состояние и одно предыдущее; движение распадается на циклы.

\begin{figure}[htbp]
\centering
\includegraphics[width=0.70\textwidth]{lecture_08_figure_12.png}
\caption{Классическая обратимая эволюция, изображённая стрелками между дискретными конфигурациями.}
\lectureframe{01:11:40}
\end{figure}

\begin{classroomqa}
\lecturetimestamp{01:11:51}
\audiencequestion{Почему, если конфигурации дискретны, шаги по времени тоже должны быть дискретными?}
\lecturerresponse{В строго классической двухсостоянийной модели нет промежуточных конфигураций, через которые могла бы проходить непрерывная траектория. Если медленно поворачивать реальную монету, её ориентация вводится как дополнительная непрерывная координата, так что это уже не модель с двумя конфигурациями. Квантовая двухсостоянийная система иная: её гильбертово пространство содержит непрерывное семейство суперпозиций, хотя ортонормированный базис имеет лишь два элемента.}
\end{classroomqa}

Состояние спина-\(\tfrac12\) может непрерывно поворачиваться от \(\ket{\uparrow}\) к \(\ket{\downarrow}\) через суперпозиции. Конечная размерность не означает конечного множества векторов состояния. Это делает возможной непрерывную временную эволюцию в дискретной квантовой системе.

Обратимая микроскопическая динамика также сохраняет различия. Два разных классических начальных условия не сливаются в одно микроскопическое состояние, потому что это стёрло бы информацию о том, какое начальное условие имело место. В гамильтоновой механике эта обратимость сопровождается теоремой Лиувилля: объём фазового пространства сохраняется.

\begin{classroomqa}
\lecturetimestamp{01:15:26}
\audiencequestion{А как насчёт нелинейных бифуркаций, при которых несколько начальных условий приближаются к одному и тому же стационарному состоянию?}
\lecturerresponse{Аттракторы возникают в эффективных диссипативных описаниях, где степени свободы окружающей среды отброшены. Фундаментальная микроскопическая гамильтонова эволюция остаётся обратимой; информация может передаваться ненаблюдаемым переменным, но различные полные состояния не превращаются буквально в одно полное состояние.}
\end{classroomqa}

Квантовая формулировка — это унитарность. Ортогональные начальные состояния переходят в ортогональные конечные состояния, и, более общо, сохраняется любое скалярное произведение.

\begin{classroomqa}
\lecturetimestamp{01:16:50}
\audiencequestion{Должны ли два состояния испытывать в точности одно и то же воздействие? Что происходит с начальными состояниями со спином вверх по \(z\) и вниз по \(z\) в магнитном поле, направленном вдоль \(x\)?}
\lecturerresponse{К обоим состояниям должен применяться один и тот же оператор эволюции. В поперечном поле оба спина прецессируют, но их векторы Блоха остаются противоположными, а векторы состояния — ортогональными. Применение разных внешних воздействий во времени означало бы сравнение разных операторов эволюции.}
\end{classroomqa}

\section{Линейность и унитарность}

Первый принцип квантовой эволюции во времени — линейность:
\begin{equation}
\ket{\psi(t)}=U(t)\ket{\psi(0)}.
\end{equation}
Оператор зависит от прошедшего интервала времени и удовлетворяет условию
\begin{equation}
U(0)=I.
\end{equation}
Как и линейность самой квантовой механики, это правило оправдывается его эмпирическим успехом.

\begin{figure}[htbp]
\centering
\includegraphics[width=0.70\textwidth]{lecture_08_figure_13.png}
\caption{Состояние в момент \(t\) получается действием линейного оператора \(U(t)\) на начальное состояние.}
\lectureframe{01:19:00}
\end{figure}

\begin{classroomqa}
\lecturetimestamp{01:19:01}
\audiencequestion{Чему должен быть равен \(U(t)\) при \(t=0\)?}
\lecturerresponse{Это должен быть единичный оператор. При нулевом интервале времени состояние остаётся неизменным для любого возможного начального вектора.}
\end{classroomqa}

Второй принцип — сохранение скалярных произведений:
\begin{equation}
\braket{\phi(t)}{\psi(t)}
=\braket{\phi(0)}{\psi(0)}.
\end{equation}
Когда \(\phi=\psi\), это сохраняет нормировку и, следовательно, полную вероятность. Когда начальные состояния ортогональны, это сохраняет идеальную различимость. Для произвольных состояний это сохраняет всё их геометрическое соотношение.

Вектор-бра эволюционирует согласно
\begin{equation}
\bra{\phi(t)}
=\bra{\phi(0)}U^\dagger(t),
\end{equation}
поэтому
\begin{align}
\braket{\phi(t)}{\psi(t)}
&=\bra{\phi(0)}U^\dagger(t)U(t)\ket{\psi(0)} \\
&=\braket{\phi(0)}{\psi(0)}.
\end{align}
Поскольку это равенство выполняется для любой пары начальных векторов,
\begin{equation}
\boxed{U^\dagger(t)U(t)=I}.
\end{equation}

\begin{figure}[htbp]
\centering
\includegraphics[width=0.74\textwidth]{lecture_08_figure_04.png}
\caption{Эволюция кета, эволюция бра и условие на скалярное произведение, приводящее к \(U^\dagger U=I\).}
\lectureframe{01:23:30}
\end{figure}

\begin{classroomqa}
\lecturetimestamp{01:22:10}
\audiencequestion{Почему бра получает \(U^\dagger\), а не \(U\)?}
\lecturerresponse{Взятие эрмитова сопряжения обращает порядок множителей и заменяет каждый оператор его эрмитово сопряжённым:
\[
\bigl(U\ket{\phi}\bigr)^\dagger
=\bra{\phi}U^\dagger.
\]
Именно такое действие в дуальном пространстве требуется для согласованности скалярных произведений.}
\end{classroomqa}

\begin{classroomqa}
\lecturetimestamp{01:24:34}
\audiencequestion{Если бы \(U\) было просто комплексным числом, что означало бы \(U^\dagger U=1\)?}
\lecturerresponse{Это означало бы \(|U|=1\), так что \(U=e^{i\theta}\) лежит на единичной окружности. Унитарная матрица — операторный аналог чистой фазы: она сохраняет длины и скалярные произведения.}
\end{classroomqa}

\section{Бесконечно малая эволюция и гамильтониан}

Пусть \(\epsilon\) — очень малый интервал времени. Поскольку \(U(0)=I\), запишем разложение до первого порядка
\begin{equation}
U(\epsilon)=I-i\epsilon h+O(\epsilon^2).
\end{equation}
Знак минус и множитель \(i\) — удобные соглашения. Взятие эрмитова сопряжения даёт
\begin{equation}
U^\dagger(\epsilon)
=I+i\epsilon h^\dagger+O(\epsilon^2).
\end{equation}
Тогда унитарность требует
\begin{align}
I
&=U^\dagger(\epsilon)U(\epsilon) \\
&=\left(I+i\epsilon h^\dagger\right)
\left(I-i\epsilon h\right)+O(\epsilon^2) \\
&=I+i\epsilon(h^\dagger-h)+O(\epsilon^2).
\end{align}
Член первого порядка должен обращаться в нуль:
\begin{equation}
\boxed{h^\dagger=h}.
\end{equation}
Генератор бесконечно малого преобразования эрмитов.

\begin{figure}[htbp]
\centering
\includegraphics[width=0.76\textwidth]{lecture_08_figure_14.png}
\caption{Унитарность в первом порядке вынуждает генератор бесконечно малого преобразования быть эрмитовым.}
\lectureframe{01:29:35}
\end{figure}

Во введённом выше временном соглашении \(h\) имеет единицы угловой частоты. Возвращая физические единицы, определим оператор энергии
\begin{equation}
H=\hbar h.
\end{equation}
Этот эрмитов оператор — гамильтониан, а его собственные значения — энергетические уровни.

Применим бесконечно малый оператор к состоянию:
\begin{equation}
\ket{\psi(\epsilon)}
=\left(I-\frac{i\epsilon}{\hbar}H\right)\ket{\psi(0)}
+O(\epsilon^2).
\end{equation}
Вычтем начальное состояние, разделим на \(\epsilon\) и перейдём к пределу:
\begin{equation}
\frac{\mathrm d}{\mathrm dt}\ket{\psi(t)}
=-\frac{i}{\hbar}H\ket{\psi(t)}.
\end{equation}
Эквивалентно,
\begin{equation}
\boxed{
i\hbar\frac{\mathrm d}{\mathrm dt}\ket{\psi(t)}
=H\ket{\psi(t)}
}.
\end{equation}
Это общее уравнение Schr\"odinger.

\begin{figure}[htbp]
\centering
\includegraphics[width=0.76\textwidth]{lecture_08_figure_15.png}
\caption{Бесконечно малое разностное отношение становится производной состояния по времени.}
\lectureframe{01:32:18}
\end{figure}

Гамильтониан зависит от физической постановки. Он может меняться со временем, если меняется внешнее магнитное поле или другой управляющий параметр. Когда постановка не зависит от времени, \(H\) является фиксированным эрмитовым оператором.

\begin{classroomqa}
\lecturetimestamp{01:34:30}
\audiencequestion{Следует ли из сохранения скалярных произведений, что системы, изначально не запутанные, должны оставаться незапутанными?}
\lecturerresponse{Нет. Унитарность сохраняет скалярные произведения между полными векторами состояния, а не структуру произведения относительно выбранного разбиения. Гамильтониан взаимодействия может перевести состояние-произведение в запутанное состояние. Если бы унитарные операторы не могли делать этого, измерения не могли бы коррелировать систему с прибором.}
\end{classroomqa}

\section{Единицы энтропии}

\begin{classroomqa}
\lecturetimestamp{01:35:08}
\audiencequestion{Имеет ли запутанность физическую единицу измерения, например джоули?}
\lecturerresponse{Энтропия запутанности — это мера информации, а не энергия. При натуральных логарифмах она измеряется в натах; при логарифмах по основанию два — в битах. И то и другое — безразмерные счётные соглашения.}
\end{classroomqa}

Термодинамические обозначения могут скрывать этот момент, потому что постоянная Больцмана переводит между температурой, измеряемой в кельвинах, и энергией:
\begin{equation}
\expect{K}_{\mathrm{monatomic}}
=\frac32 k_{\mathrm B}T.
\end{equation}
Для обратимого теплообмена справедливо соотношение
\begin{equation}
\delta Q_{\mathrm{rev}}=T\,\mathrm dS_{\mathrm{th}}.
\end{equation}
Если информационная энтропия равна
\begin{equation}
S_{\mathrm{info}}=-\sum_i p_i\log p_i,
\end{equation}
то обычная термодинамическая энтропия равна
\begin{equation}
S_{\mathrm{th}}=k_{\mathrm B}S_{\mathrm{info}}.
\end{equation}
Полагая \(k_{\mathrm B}=1\), мы измеряем температуру в единицах энергии и оставляем энтропию безразмерной. Тогда постоянная Больцмана понимается как коэффициент перевода, а не как новый вид единицы запутанности. Историческое отступление, связывающее статистические идеи Больцмана с работой Эйнштейна 1905 года о броуновском движении, не изменяет информационно-теоретического определения.

\section{Энергетические собственные состояния и фаза}

\begin{classroomqa}
\lecturetimestamp{01:39:54}
\audiencequestion{Почему эрмитов генератор \(H\) отождествляется с энергией?}
\lecturerresponse{Сама по себе унитарность доказывает лишь, что генератор эрмитов. Его отождествление с энергией — дополнительное физическое утверждение: энергия порождает трансляции по времени. Для не зависящего от времени гамильтониана тот же оператор, который управляет эволюцией, сохраняется, а его собственные значения являются стационарными энергетическими уровнями. Конкретные гамильтонианы и придают этому отождествлению экспериментальное содержание.}
\end{classroomqa}

Связь между постоянными Планка и частотой имеет вид
\begin{equation}
E=h\nu=\hbar\omega,
\qquad
\omega=2\pi\nu,
\qquad
h=2\pi\hbar.
\end{equation}

\begin{classroomqa}
\lecturetimestamp{01:40:30}
\audiencequestion{Должно ли уравнение эволюции содержать \(h\), \(\hbar\) или дополнительный множитель \(2\pi\)?}
\lecturerresponse{Используя угловую частоту \(\omega\), получаем уравнение \(E=\hbar\omega\), и уравнение Schr\"odinger содержит \(\hbar\). Используя обычную частоту \(\nu\), получаем эквивалентное соотношение \(E=h\nu\). Множитель \(2\pi\) полностью учитывается соотношениями \(\omega=2\pi\nu\) и \(h=2\pi\hbar\).}
\end{classroomqa}

Предположим, что начальное состояние является энергетическим собственным вектором:
\begin{equation}
H\ket{\psi(0)}=E\ket{\psi(0)}.
\end{equation}

\begin{classroomqa}
\lecturetimestamp{01:43:18}
\audiencequestion{Если состояние остаётся энергетическим собственным вектором, означает ли это, что оно не меняется?}
\lecturerresponse{Оно может приобрести зависящий от времени комплексный множитель и остаться собственным вектором с тем же собственным значением. Унитарность ограничивает этот множитель фазой. Таким образом, луч стационарен, хотя кет вращается по фазе.}
\end{classroomqa}

Запишем
\begin{equation}
\ket{\psi(t)}=f(t)\ket{\psi(0)}.
\end{equation}

\begin{figure}[H]
\centering
\includegraphics[width=0.70\textwidth]{lecture_08_figure_05.png}
\caption{Энергетическое собственное состояние остаётся в том же направлении в пространстве Гильберта с точностью до скалярного множителя \(f(t)\).}
\lectureframe{01:43:54}
\end{figure}

Подстановка в уравнение Schr\"odinger даёт
\begin{align}
i\hbar\frac{\mathrm d f}{\mathrm dt}\ket{\psi(0)}
&=f(t)H\ket{\psi(0)} \\
&=E f(t)\ket{\psi(0)}.
\end{align}
Сократим общий вектор:
\begin{equation}
\frac{\mathrm d f}{\mathrm dt}
=-\frac{iE}{\hbar}f(t).
\end{equation}
При \(f(0)=1\),
\begin{equation}
f(t)=e^{-iEt/\hbar},
\qquad
\ket{\psi(t)}=e^{-iEt/\hbar}\ket{\psi(0)}.
\end{equation}

\begin{figure}[H]
\centering
\includegraphics[width=0.76\textwidth]{lecture_08_figure_16.png}
\caption{Уравнение на собственные значения энергии и возникающая из него фазовая эволюция.}
\lectureframe{01:46:10}
\end{figure}

\begin{classroomqa}
\lecturetimestamp{01:45:39}
\audiencequestion{Каков коэффициент при \(t\) в фазе \(e^{-iEt/\hbar}\)?}
\lecturerresponse{Это угловая частота,
\[
\omega=\frac{E}{\hbar}.
\]
Эквивалентно, \(E=\hbar\omega\). Вещественный коэффициент в показателе экспоненты описывал бы скорость роста или затухания; мнимый показатель здесь описывает унитарное вращение фазы.}
\end{classroomqa}

Это соотношение между фазой и частотой и уравнение Schr\"odinger — две формы одного и того же закона эволюции во времени. Это уравнение является общим; свободная частица, частица в потенциале и спин в магнитном поле различаются своими гамильтонианами. Выбор этих гамильтонианов и их решение — следующий шаг.
